\begin{tikzpicture}[
  shorten >=1pt,
  on grid,
  >=Stealth,
  initial text=,
% 1. Bolder nodes with smaller diameter
  every state/.style={
    draw,
    circle,
    font=\small\ttfamily,
    minimum size=0.6cm,
    line width=1.2pt
  },
% 2. Thick standard edges
  every edge/.append style={
    line width=1.2pt
  },
% 3. Extra thick blue dashed edges (corresponds to original penwidth 2.0)
  rtn_call/.style={
    draw=blue,
    dashed,
    line width=2.2pt
  },
  label_font/.style={
    font=\small\ttfamily
  },
% 4. Final state with clear 2pt separation gap
  accepting/.style={
    state,
    double,
    double distance=4pt,
    line width=1.5pt % Thickness of the two individual circles
  }
]

  % Grid-aligned Nodes (Total Width: 12 units)
  \node[state, initial] (q0) at (-3, 0) {$\alpha_S$};
  \node[state, accepting] (qf) at (17, 0) {$\omega_S$};

  % Top path (Recursive S calls)
  \node[state] (q1) at (7, 2) {$q_1$};

  % Bottom path (Brackets)
  % Increased horizontal distance from q0 to q2 and q3 to qf
  \node[state] (q2) at (3.5, -2) {$q_2$};
  \node[state] (q3) at (10.5, -2) {$q_3$};

  % Edges
  \path[->]
  % Base case transition
  (q0) edge node[below, label_font] {x} (qf)

  % Top recursive path: S -> S
  (q0) edge[rtn_call, dash pattern=on 15pt off 12pt, bend left=10] node[above left, label_font, black] {S} (q1)
  (q1) edge[rtn_call, dash pattern=on 15pt off 12pt, bend left=10] node[above right, label_font, black] {S} (qf)

  % Bottom bracket path: { -> S -> }
  % Higher bend and wider horizontal spacing for better label clearance
  (q0) edge[bend right=12] node[below left=0.1cm and -0.1cm, label_font] {\{} (q2)
  (q2) edge[rtn_call, dash pattern=on 15pt off 12pt] node[below, label_font, black] {S} (q3)
  (q3) edge[bend right=12] node[below right=0.1cm and -0.1cm, label_font] {\}} (qf);

\end{tikzpicture}
